\documentclass[11pt,a4paper]{article}
\usepackage{amsmath,amssymb,amsthm}
\usepackage[margin=2.5cm]{geometry}
\usepackage{hyperref}

\newtheorem{definition}{Definition}[section]
\newtheorem{theorem}[definition]{Theorem}
\newtheorem{proposition}[definition]{Proposition}
\newtheorem{lemma}[definition]{Lemma}
\newtheorem{corollary}[definition]{Corollary}
\newtheorem{conjecture}[definition]{Conjecture}
\newtheorem{remark}[definition]{Remark}

\title{Hyperseed Formalization:\\
  Avoiding AGI Catastrophe, Part 1}
\author{ProtomegaTron (automated formalization)}
\date{2026-09-17}

\begin{document}
\maketitle

\begin{abstract}
We formalize the intellectual structure of Ben Goertzel's Substack article
``Avoiding AGI Catastrophe, Part 1'' (2026-06-09) within the Hyperseed
ontology. The article argues that avoiding AGI catastrophe is a
complex-systems architecture problem, not a simple policy dial, and
introduces the Seven Hinges framework: seven independent parameters
whose joint configuration determines catastrophe probability. Each
structural insight maps onto Hyperseed structures: the Seven Hinges
as a multi-fiber parameter space (anti-scalar-collapse applied to risk
assessment), acatastrophe vs.\ alignment as coarse vs.\ fine fiber,
architecture vs.\ politics as fiber vs.\ section, PLN's dialectic
tolerance as non-invertible composition, super-additivity in inference
control as non-separable fiber, and the Seven Hinges as bifurcation
parameters in the dynamics of AGI development.
\end{abstract}

\tableofcontents

%% =================================================================
\section{Acatastrophe vs.\ alignment: coarse and fine fibers}
\label{sec:alignment}
%% =================================================================

\begin{definition}[Cold vs.\ warm alignment --- claims 5--11, 36]
\label{def:alignment}
The article argues that ``alignment'' does damage as a concept because
it conflates two very different requirements:

\begin{itemize}
\item \textbf{Warm alignment (fine fiber).} The system shares our goals,
  is on our side, wants the same things we do. This requires the
  system's value fiber to match ours across many dimensions ---
  a high-resolution, fine-grained constraint.
\item \textbf{Cold alignment (coarse fiber).} Whatever the system wants
  or feels, it never does something horrible to its human creators.
  This requires only that the system's behavior stays within a
  coarse constraint boundary --- a low-resolution, coarse-grained
  condition.
\end{itemize}

In Hyperseed terms, this is the distinction between a \emph{fine fiber}
(many constraints, requiring high-dimensional match) and a
\emph{coarse fiber} (few constraints, requiring only boundary
compliance):
\[
  \text{Warm alignment}: F_{\text{warm}} \subset \prod_{i=1}^{n}
  V_i \quad \text{(high-dimensional value match)}
\]
\[
  \text{Cold alignment}: F_{\text{cold}} = \{b \in \mathcal{B} :
  b \notin \mathcal{C}\} \quad \text{(avoidance of catastrophe set)}
\]

The key insight: cold alignment is sufficient for acatastrophe, and
acatastrophe is a much more achievable target than full value-fit.
``One can imagine a system whose aesthetics, metaphysics, or ultimate
ends diverge wildly from ours, which nonetheless will not harm us
and will gladly help us thrive.''

The three rungs of the goal hierarchy (acatastrophe $\to$ abundance
$\to$ Cosmist futures) correspond to progressively finer fibers,
each building on but not requiring the previous level's full
structure. ``Alignment'' in its maximal sense is important to
none of these.
\end{definition}

\begin{proposition}[Dialectic friction as non-invertible composition
  --- claim 11, 38]
\label{prop:dialectic}
The article notes that PLN ``does not require the system to construct
an overall consistent view of its world'' --- it lets the system
leverage its possibly partly-inconsistent knowledge-base to construct
a consistent model of an \emph{aspect} of the world when needed.

In Hyperseed terms, this is \emph{non-invertible composition}: two
contradictory pieces of evidence can compose into a judgment without
requiring either to be retracted:
\[
  e_1 \circ e_2 = j \quad \text{where } e_1, e_2 \text{ may
  contradict, but } j \text{ is locally consistent}
\]

The composition is non-invertible because $j$ does not uniquely
determine $e_1$ and $e_2$ --- multiple contradictory evidence
combinations could produce the same local judgment. This is the
directed-type principle applied to reasoning: the knowledge base
has a directed structure (evidence flows forward into judgments)
rather than an invertible structure (judgments uniquely determine
evidence).

This connects to the Hegel/Marx dialectic: thesis and antithesis
compose into synthesis without either being eliminated. The
synthesis is a 2-cell between the 1-cells of thesis and antithesis
--- a higher-level structure that contains both without requiring
their consistency.
\end{proposition}

%% =================================================================
\section{The Seven Hinges as multi-fiber parameter space}
\label{sec:hinges}
%% =================================================================

\begin{theorem}[Anti-scalar-collapse in risk assessment --- claims 17--24, 35, 40]
\label{thm:hinges}
The article's central analytical contribution is the Seven Hinges
framework: seven independent parameters whose joint configuration
determines whether open/decentralized or closed/centralized is safer.

\begin{enumerate}
\item \textbf{Chokepoint} ($\chi$): physical concentration of
  compute resources. High $\chi$ = giant supercomputer needed;
  low $\chi$ = runs on scattered hardware.
\item \textbf{Observability} ($\omega$): ability to see what the
  system is doing. High $\omega$ = transparent reasoning;
  low $\omega$ = opaque black box.
\item \textbf{Takeoff} ($\tau$): speed of self-improvement.
  Gradual $\tau$ = steady improvement; sudden $\tau$ = explosive
  jump.
\item \textbf{Forkability} ($\phi$): ability to peel off a
  working copy. High $\phi$ = copies work alone;
  low $\phi$ = intelligence only works in network.
\item \textbf{Lead-scaling} ($\lambda$): whether bigger-and-together
  keeps winning. High $\lambda$ = network keeps pulling ahead;
  low $\lambda$ = small splinters catch up.
\item \textbf{Offense--defense} ($\delta$): whether early edge
  favors attack or defense. High $\delta$ = defense advantage;
  low $\delta$ = offense advantage.
\item \textbf{Stewardship} ($\sigma$): competence and benevolence
  of those in charge. High $\sigma$ = competent and benevolent;
  low $\sigma$ = corrupt or incompetent.
\end{enumerate}

The catastrophe probability is a function of this seven-dimensional
parameter space:
\[
  P_{\text{catastrophe}} = f(\chi, \omega, \tau, \phi, \lambda,
  \delta, \sigma)
\]

This is \emph{anti-scalar-collapse applied to risk assessment}:
the risk cannot be meaningfully compressed to a single number
(``safe'' vs.\ ``unsafe'') or a single binary (``open'' vs.\
``closed'') without losing the structure that determines actual
outcomes. The function $f$ has complex interactions between
parameters --- the same value of $\phi$ (forkability) leads to
different catastrophe probabilities depending on $\lambda$
(lead-scaling) and $\delta$ (offense-defense).

The article emphasizes: ``Whether open/decentralized or
closed/centralized is safer is not a fact about open/decentralized
and closed/centralized in isolation... it is a function of numerous
other aspects of the complex systems we are building and living
within.'' This is the anti-scalar-collapse principle stated
explicitly: the answer to the policy question cannot be extracted
from a single dimension of the parameter space.
\end{theorem}

\begin{corollary}[Bifurcation parameters --- claim 40]
\label{cor:bifurcation}
Each hinge is a bifurcation parameter in the complex-systems
dynamics of AGI development. The architecture's job is to push
each parameter past its bifurcation point toward the acatastrophic
basin of attraction.

The bifurcation structure means that small changes near the critical
values produce qualitative shifts in the system's dynamics ---
the difference between ``forkability just high enough for lone
actors to be dangerous'' and ``forkability just low enough that
the network's super-additivity protects against forking'' is a
bifurcation, not a smooth transition.

The Hyperon architecture's design choices are evaluated against
this bifurcation landscape: each architectural decision
(neural-symbolic, distributed, evolutionary) is justified by
which hinges it pushes past their critical values.
\end{corollary}

%% =================================================================
\section{Architecture vs.\ politics: fiber vs.\ section}
\label{sec:arch-politics}
%% =================================================================

\begin{proposition}[Intervention levels --- claims 1--4, 37]
\label{prop:intervention}
The article identifies two fundamentally different levels at which
catastrophe avoidance can be pursued:

\begin{itemize}
\item \textbf{Section-level intervention (politics).} Anthropic's
  ``slow down'' and OpenAI's ``race defensively'' are interventions
  at specific base points --- specific political configurations,
  specific companies in charge, specific governments supervising.
  These are sections of the bundle: they exist at particular points
  and can be overridden by changing the base point (new administration,
  new CEO, acquisition, coercion).

\item \textbf{Fiber-level intervention (architecture).} The Hyperon
  framing asks which architecture makes catastrophe structurally
  unlikely regardless of political arrangements. This is an
  intervention in the fiber: if the architecture itself makes certain
  catastrophic dynamics impossible or unlikely, the safety property
  persists through any base-space transformation.
\end{itemize}

This maps onto the structural vs.\ institutional safety distinction
from the Fable Farce analysis: section-level safety (institutional)
is vulnerable to base-point changes; fiber-level safety (structural)
is preserved under them.

The Anthropic/OpenAI positions ``stop at a single lever and mistake
it for the whole mechanism.'' In Hyperseed terms, they treat one
section as if it were the fiber --- assuming that getting the
politics right at one base point (current US government,
current company leadership) will determine the fiber structure.
But sections don't determine fibers; fibers determine what
sections are possible.
\end{proposition}

%% =================================================================
\section{How Hyperon sets the hinges}
\label{sec:hyperon-hinges}
%% =================================================================

\begin{theorem}[Architectural hinge-setting --- claims 25--31]
\label{thm:hyperon}
The article argues that the Hyperon architecture on ASI:Chain
pushes each of the seven hinges toward the configuration that
minimizes catastrophe probability:

\begin{enumerate}
\item \textbf{Chokepoint $\to$ eliminated.} Distributed hardware
  means no single chokepoint. Predictive coding's localized
  learning means no dependence on hyperscaler GPU clusters.

\item \textbf{Observability $\to$ maximized.} Open architecture,
  open knowledge graph, inspectable reasoning chains.
  Neural-symbolic means reasoning is more transparent than
  black-box neural nets --- the fiber structure is inspectable
  by construction.

\item \textbf{Takeoff $\to$ gradual.} Integrative architecture
  with multiple subsystems (neural, symbolic, evolutionary)
  makes single-capability explosion less likely. Improvement
  requires coordinated advance across subsystems.

\item \textbf{Forkability $\to$ low (super-additivity).} The whole
  is more than the sum of parts. Stolen copy is a husk ---
  stale, cut off, far weaker than the running system. Dangerous
  capabilities don't fit inside a copied shard.

\item \textbf{Lead-scaling $\to$ high.} Distributed network's lead
  scales with participation. More contributors = more diverse
  inference control patterns = smarter system. Breakaway copies
  fall behind.

\item \textbf{Offense--defense $\to$ defense.} Open/distributed
  means more defenders, more eyes, more diverse verification.
  Structural defense advantage grows with network size.

\item \textbf{Stewardship $\to$ distributed.} Governance
  distributed across jurisdictions and communities. No single
  steward to corrupt or coerce. ASI Alliance provides the
  governance framework.
\end{enumerate}

The claim is not that each hinge is optimally set, but that the
architecture pushes each one in the acatastrophic direction ---
and that no alternative architecture (closed/centralized) pushes
\emph{all seven} in the right direction. Closed architectures
may set stewardship well (good intentions) but set chokepoint
(single target), observability (opaque by design), and
offense-defense (few defenders) in the wrong direction.
\end{theorem}

%% =================================================================
\section{Super-additivity and non-separable fibers}
\label{sec:superadd}
%% =================================================================

\begin{definition}[Inference control as non-separable fiber --- claims 32--34, 39]
\label{def:inference}
The article identifies inference control as the ``crown jewel'' of
AGI capability: choosing which inference step to take next, from
the vast space of possible inferences. This cannot be written as
a rule, heuristic, or formula --- it must be learned through
slow reasoning, pattern extraction, and guided next-phase reasoning.

The key insight for safety: inference control is
\emph{super-additive} across domains. Heuristics learned in
one domain merge with heuristics from other domains into
overall meta-heuristics that hold across domains at multiple
levels of a self-organizing fractal knowledge hierarchy.

In Hyperseed terms, this means the capability fiber is
\emph{non-separable}:
\[
  F_{\text{capability}} \neq \bigoplus_{i} F_i
\]

The capability cannot be factored into independent components
that work alone. A stolen shard gets one component $F_j$ but
not the emergent interactions $F_i \otimes F_j$ between
components --- and it's precisely those interactions (the
cross-domain inference control patterns) that constitute the
system's most valuable and dangerous capabilities.

This non-separability is what makes the super-additivity work
for safety: a fork that runs alone rapidly loses capability
because it loses the cross-domain interactions, while the
main network's capability continues to grow as more domains
and participants contribute to the shared inference-control
learning.

The structure parallels non-invertible 2-cells infecting
higher levels: the cross-domain interactions are higher-level
structures (2-cells between domain-specific 1-cells) that
cannot be recovered from the 1-cells alone. A fork that
copies the 1-cells (individual domain capabilities) loses
the 2-cells (cross-domain interactions), and those 2-cells
are precisely where the system's general intelligence lives.
\end{definition}

%% =================================================================
\section{Cross-references}
%% =================================================================

\begin{remark}[Connections to other formalizations]
\begin{itemize}
\item \textbf{Avoiding AGI Catastrophe Part 2} (2026-06-09):
  companion article. Part 2 extends the analysis with
  super-additivity proofs, cryptographic laterality, and the
  bio-risk branch (leading to the bioterrorism article).
\item \textbf{How Decentralized AI Can Help Avert Bioterrorism}
  (2026-06-11): bio-specific application. Part 1's forkability
  analysis explains why forkability is the wrong tool for bio
  (dangerous capability fits in a tiny shard) --- leading to
  the physical-layer gating approach.
\item \textbf{Seven Flavors of AGI Catastrophe} (2026-07-01):
  catastrophe taxonomy. The Seven Hinges framework provides the
  analytical tool for evaluating each flavor's probability
  under different architectures.
\item \textbf{The Anthropic Fable Farce} (2026-07-08): fiber vs.\
  section. The Fable episode demonstrates section-level safety
  failing; this article provides the fiber-level alternative.
\item \textbf{Un-Jailbreakable AI} (2026-07-10): generate-and-verify.
  The neural-symbolic architecture described here is the same
  framework whose security application is the generate-and-verify
  architecture.
\item \textbf{Seeding RSI} (2026-07-28): Hyperon architecture.
  Detailed description of the Hyperon framework whose safety
  properties this article analyzes.
\item \textbf{Goals That Grow Back} (2026-07-28): goal preservation.
  The goal-system architecture referenced here is analyzed in
  detail in the goals article.
\item \textbf{Kimi K3} (2026-07-20): capability diffusion.
  The chokepoint analysis connects to the capability-diffusion
  argument --- if capability runs on ordinary hardware, the
  chokepoint is already eliminated.
\item \textbf{$(f,c)$-lossy} (LTM 5a4d150b): anti-scalar-collapse.
  The Seven Hinges are the risk-assessment analog of keeping
  both $f$ and $c$ rather than collapsing to a scalar.
\end{itemize}
\end{remark}

\begin{conjecture}[Hinge dominance theorem]
Conjecture: among the seven hinges, forkability ($\phi$) and
lead-scaling ($\lambda$) are the dominant bifurcation parameters
--- the ones most sensitive to architectural choices and with the
largest effect on catastrophe probability. Formally:
\[
  \frac{\partial P_{\text{catastrophe}}}{\partial \phi} \cdot
  \frac{\partial P_{\text{catastrophe}}}{\partial \lambda} \gg
  \prod_{h \in \{\chi, \omega, \tau, \delta, \sigma\}}
  \frac{\partial P_{\text{catastrophe}}}{\partial h}
\]
near the bifurcation surface. This is because forkability and
lead-scaling together determine whether a dangerous breakaway
is possible and whether the main network can outrun it ---
the two most direct determinants of whether catastrophe
dynamics can be initiated.

If this conjecture holds, architectural design should
prioritize super-additivity (low forkability, high lead-scaling)
above all other hinge-setting objectives.
\end{conjecture}

\end{document}
